\documentclass[12pt,a4paper]{article}
\usepackage[margin=2cm]{geometry}
\usepackage[spanish]{babel}
\usepackage{graphicx, multicol, latexsym, amsmath, amssymb}
\usepackage{array}
\usepackage{tabularx}
\usepackage{setspace}
\usepackage{titlesec}
\usepackage{subcaption}
\usepackage{float}
\usepackage{url}
\usepackage{hyperref}

% Ajuste de espacio entre secciones
\titlespacing*{\section}{0pt}{0.5em}{0.5em}

\renewcommand{\arraystretch}{1.3}

\begin{document}

% ---------------------------
% Encabezado con datos
% ---------------------------
\noindent
\begin{tabularx}{\textwidth}{|X|X|}
\hline
\textbf{Carné:} 2020045294 & \textbf{Nombre:} Justin Jaffeth Corea Masís \\
\hline
\textbf{Correo:} coreajustin288@estudiantec.cr & \textbf{Semana:} 5 del 3/09/2025 al 10/09/2025 \\
\hline
\end{tabularx}

\vspace{0.5cm}

% ---------------------------
% Tabla de actividades
% ---------------------------
\section*{Actividades realizadas}

% \begin{tabularx}{\textwidth}{|>{\raggedright\arraybackslash}p{8cm}|c|c|}
\begin{tabularx}{\textwidth}{|>{\raggedright\arraybackslash}p{12cm}|c|c|}
\hline
\textbf{Actividad} & \textbf{Fecha} & \textbf{Horas} \\
\hline
	Trabajé el documento escrito dejando lugares imágenes y corrigiendo cosas que cambié. & 3/09/2025 & 3 \\
\hline
	Diseño de experimentos para mejorar YOLO y el OCR. Dejar código listo para ejecutarlos & 6/09/2025 & 4 \\
\hline
	Arreglar problemas con el clúster Kabré porque se desconfiguró & 7/09/2025 & 8 \\
\hline
	Ejecución de experimentos YOLO y OCR & 8/09/2025 & 8 \\
\hline 
	Diseñé más experimentos para el OCR y los ejecuté & 9/09/2025 & 4 \\
\hline
	Documentación y trabajo escrito & 9/09/2025 & 4 \\
\hline
\end{tabularx}

\vspace{1cm}

% ---------------------------
% Firma
% ---------------------------
\noindent
\textbf{Firma del Profesor Asesor:}

\vspace{2cm}

\noindent\rule{8cm}{0.4pt}  
Nombre y firma

\newpage

\section{Detección de placas}
Realicé los mismos experimentos previos, pero agregando ahora imágenes de vehículos
internacionales tanto en el conjunto de entrenamiento como en validación.

Los resultados de entrenamiento usando imágenes de Costa Rica e internaciones mezcladas
sigue dando mejores resultados que cuando se usan solo internacionales o solo nacionales
para entrenar.

\section{Reconocimiento de placas}

En los primeros experimentos 1-4 usé el mismo aumentado de datos y cambié
solo los parámetros de imagen, optimizador y red.
De los experimentos 5-8 repetí los parámetros de experimentos 1-4, pero cambié
el aumentado de datos.
Intenté variar el tipo de aumentado de datos en los experimentos 5-8.

Aumentado base: 
Combinaba gran cantidad de transformaciones con gran intensidad:
\begin{figure}[H]
	\begin{subfigure}[b]{0.3\textwidth}
		\includegraphics[width=0.5\textwidth]{../../src/sources/images/images/synthetic_plate_5bcc130c.png}
	\end{subfigure}
	\begin{subfigure}[b]{0.3\textwidth}
		\includegraphics[width=0.5\textwidth]{../../src/sources/images/images/synthetic_plate_48752a1d.png}
	\end{subfigure}
	\begin{subfigure}[b]{0.3\textwidth}
		\includegraphics[width=0.5\textwidth]{../../src/sources/images/images/synthetic_plate_2c06fd53.png}
	\end{subfigure}

	\begin{subfigure}[b]{0.3\textwidth}
		\includegraphics[width=0.5\textwidth]{../../src/sources/images/images/synthetic_plate_fa402bc9.png}
	\end{subfigure}
	\begin{subfigure}[b]{0.3\textwidth}
		\includegraphics[width=0.5\textwidth]{../../src/sources/images/images/synthetic_plate_98be42f6.png}
	\end{subfigure}
	\begin{subfigure}[b]{0.3\textwidth}
		\includegraphics[width=0.5\textwidth]{../../src/sources/images/images/synthetic_plate_8eb9619e.png}
	\end{subfigure}
\end{figure}

Experimento 5: Aplica ligero cambio en perspectiva y pequeños obstáculos visuales.

\begin{figure}[H]
	\includegraphics[width=0.5\textwidth]{/home/akai/Pictures/Screenshots/augmentation-1.png}
\end{figure}

Experimento 4: Agrega un poco de brillo y ruido gausiano.

\begin{figure}[H]
	\includegraphics[width=0.5\textwidth]{/home/akai/Pictures/Screenshots/augmentation-2.png}
\end{figure}

Experimento 6: Combina los casos previos y aumenta muy ligeramente los dos previos casos.

\begin{figure}[H]
	\includegraphics[width=0.5\textwidth]{/home/akai/Pictures/Screenshots/augmentation-3.png}
\end{figure}

Experimento 7: Se agrega pixelado y se aumenta el desenfoque.

\begin{figure}[H]
	\includegraphics[width=0.5\textwidth]{/home/akai/Pictures/Screenshots/augmentation-4.png}
\end{figure}

Experimento 8: Se combinan los casos previos y se aumentan ligeramente de nuevo.

\begin{figure}[H]
	\includegraphics[width=0.5\textwidth]{/home/akai/Pictures/Screenshots/augmentation-5.png}
\end{figure}

Resultados de entrenamiento y validación:

\begin{figure}[H]
	\begin{subfigure}[b]{0.9\textwidth}
		\includegraphics[width=0.9\textwidth]{/home/akai/Pictures/Screenshots/train-experimentos.png}
		\caption{Entrenamiento}
	\end{subfigure}

	\begin{subfigure}[b]{0.9\textwidth}
		\includegraphics[width=0.9\textwidth]{/home/akai/Pictures/Screenshots/val-experimentos.png}
		\caption{Validación}
	\end{subfigure}
\end{figure}


\begin{figure}[H]
	\begin{subfigure}[b]{0.9\textwidth}
		\includegraphics[width=0.8\textwidth]{/home/akai/Pictures/Screenshots/imagen.png}
		\caption{Parámetros de imagen}
	\end{subfigure}

	\begin{subfigure}[b]{0.9\textwidth}
		\includegraphics[width=0.8\textwidth]{/home/akai/Pictures/Screenshots/optimizador.png}
		\caption{Parámetros de optimizador}
	\end{subfigure}

	\begin{subfigure}[b]{0.9\textwidth}
		\includegraphics[width=0.8\textwidth]{/home/akai/Pictures/Screenshots/modelo.png}
		\caption{Parámetros del modelo}
	\end{subfigure}
\end{figure}

\end{document}
